\section{Testbed}
\label{sec:testbed}

The testbed is designed so that the deployment topology is the only thing that differs between
configurations. This section describes the application, the four configurations, the resource
budget and the per-layer latency accounting; Section~\ref{sec:method} describes how it was exercised.

\subsection{One artifact, four topologies}
\label{sec:artifact}

The application is a small e-commerce back end written in Java~25 with Spring Boot~4.1.1 (Tomcat, JDBC,
HikariCP) and PostgreSQL~17. Its business logic is split into eight modules, each owning its tables:
\emph{Auth} (sessions), \emph{Profiles} (users and loyalty tier), \emph{Search} (a category index),
\emph{Products} (catalogue), \emph{Stock} (inventory levels), \emph{Reservations}, \emph{Carts} and
\emph{Orders}. Each module is a Java interface with one implementation. All eight modules are compiled into a
single artifact, and a process reads one setting, \texttt{SERVICE\_MODULES}, that lists the modules it hosts.
A call to a hosted module is a plain method call. A call to any other module goes through a generic stub that
serialises the arguments to JSON, sends \texttt{POST /rpc/\{Module\}/\{method\}} over HTTP/1.1 with
keep-alive connections, and deserialises the result. Because every configuration runs the same binary with
the same business code, a difference between configurations cannot come from a difference in implementation.
This is a stronger form of the shared-library approach described in our earlier protocol, where equivalence
had to be enforced by tests.

\begin{table}[t]
\caption{Configurations. All share one artifact; only the module-to-process mapping differs.
CPU is the per-pod share of the equal budget (\xo{appcpu} cores for the application tier and
\xo{dbcpu} cores for the database tier).}
\label{tab:configs}
\centering\small
\begin{tabular}{@{}llrrr@{}}
\toprule
Config & Services (modules) & Pods & CPU/app pod & CPU/db pod \\
\midrule
A & \texttt{mono} (all eight) & 1+1 & 4.0 & 2.0 \\
B & \texttt{user} (Auth, Profiles), \texttt{catalog} (Search, Products), & 4+4 & 1.0 & 0.5 \\
  & \texttt{inventory} (Stock, Reservations), \texttt{order} (Carts, Orders) & & & \\
C & one service per module & 8+8 & 0.5 & 0.25 \\
D & as C, every pod with an Istio sidecar & 8+8 & 0.5 (+ sidecar) & 0.25 (+ sidecar) \\
\bottomrule
\end{tabular}
\end{table}

Table~\ref{tab:configs} lists the configurations. Config~A is the monolith. Config~B groups the modules into
four services along domain lines, the grouping most teams would choose first. Config~C gives every module its own
service. Config~D is C with an Istio~1.31 service mesh injected into every pod, including the database pods;
it isolates the mesh as the single added variable. The earlier protocol's Config~D also added asynchronous
messaging, which would have changed two variables at once; we dropped it for that reason. Every service has
its own PostgreSQL instance holding only the tables of its modules, so Config~A has one database and
Configs~C and~D have eight. In~D, Istio runs with its default PERMISSIVE peer authentication and automatic
mutual TLS, so every sidecar-to-sidecar connection, including those to databases, is encrypted; only the
load generator's connection to the edge service is plaintext.

\subsection{Operations and call graphs}
\label{sec:ops}

Two edge operations exercise the system. \emph{Browse} is a read-only listing: it verifies the session,
reads twenty product identifiers for a category page from the Search index, fetches those products and
their stock levels, and returns about 5.5\,kB of JSON. \emph{Checkout} is a deep write path: it verifies the
session, reads the user's loyalty tier and cart, prices the cart, reserves stock (Reservations calls Stock,
which decrements all items in one atomic \texttt{UPDATE}) and records the order. Every cart contains one of
ten popular products, so concurrent checkouts contend for a small set of hot rows, the kind of shared-state
coherency cost the USL's $\beta$ term is meant to capture. The data set has 10{,}000 users, 10{,}000
products in 50 categories and three cart items per user.

\begin{table}[t]
\caption{Remote calls per request. Each call is one JSON/HTTP round trip between processes; the counts were
confirmed at run time from the \texttt{calls} field every response carries.}
\label{tab:hops}
\centering\small
\begin{tabular}{@{}lrrrrr@{}}
\toprule
Operation & DB statements & A & B & C & D \\
\midrule
Browse   & 4 & 0 & 2 & 3 & 3 \\
Checkout & 7 & 0 & 4 & 6 (depth 2) & 6 (depth 2) \\
\bottomrule
\end{tabular}
\end{table}

\begin{figure}[t]
\centering
\begin{tikzpicture}[
    font=\footnotesize,
    mod/.style={draw, rounded corners=2pt, minimum width=19mm, minimum height=6.5mm, fill=black!4, align=center},
    edge/.style={draw, rounded corners=2pt, minimum width=19mm, minimum height=6.5mm, fill=black!12, align=center, thick},
    grp/.style={draw, dashed, rounded corners=4pt, inner sep=2.2mm, black!55},
    br/.style={-{Stealth[length=1.6mm]}, thick, xoBrowse},
    co/.style={-{Stealth[length=1.6mm]}, thick, dashed, xoCheckout},
    lbl/.style={font=\scriptsize\itshape, black!65}
]
\node[mod]  (auth)   at (0, 0)     {Auth};
\node[mod]  (prof)   at (0, -1.5)  {Profiles};
\node[edge] (search) at (3.1, 0)   {Search};
\node[mod]  (prod)   at (3.1, -1.5){Products};
\node[mod]  (stock)  at (6.2, 0)   {Stock};
\node[mod]  (resv)   at (6.2, -1.5){Reservations};
\node[edge] (orders) at (9.3, 0)   {Orders};
\node[mod]  (carts)  at (9.3, -1.5){Carts};

\node[grp, fit=(auth)(prof)]    (g1) {};
\node[grp, fit=(search)(prod)]  (g2) {};
\node[grp, fit=(stock)(resv)]   (g3) {};
\node[grp, fit=(orders)(carts)] (g4) {};
\node[lbl, below=0.5mm of g1] {\texttt{user}};
\node[lbl, below=0.5mm of g2] {\texttt{catalog}};
\node[lbl, below=0.5mm of g3] {\texttt{inventory}};
\node[lbl, below=0.5mm of g4] {\texttt{order}};

% browse (solid): Search -> Auth, Products, Stock
\draw[br] (search.west) -- (auth.east);
\draw[br] (search.south) -- (prod.north);
\draw[br] (search.east) -- (stock.west);
% checkout (dashed): Orders -> Auth, Profiles, Carts, Products, Reservations; Reservations -> Stock
\draw[co] (orders.north) .. controls (8.2, 1.35) and (1.2, 1.35) .. (auth.north);
\draw[co] (orders.south west) .. controls (7.6, -3.35) and (1.4, -3.35) .. (prof.south east);
\draw[co] (orders.south) -- (carts.north);
\draw[co] (orders.south west) .. controls (7.6, -0.75) and (4.6, -0.75) .. (prod.north east);
\draw[co] (orders.south west) -- (resv.north east);
\draw[co] (resv.north) -- (stock.south);

\draw[-{Stealth[length=1.6mm]}, thick] (3.1, 1.15) node[above, font=\scriptsize] {client: browse} -- (search.north);
\draw[-{Stealth[length=1.6mm]}, thick] (10.7, 0.9) node[above, font=\scriptsize] {client: checkout} -- (orders.north east);

\draw[br] (0.2, -3.95) -- ++(0.7, 0) node[right, black, font=\scriptsize] {browse call};
\draw[co] (3.2, -3.95) -- ++(0.7, 0) node[right, black, font=\scriptsize] {checkout call};
\node[draw, dashed, rounded corners=3pt, black!55, minimum width=6mm, minimum height=3mm] at (6.55, -3.95) {};
\node[right, font=\scriptsize] at (6.9, -3.95) {service in Config~B};
\end{tikzpicture}
\caption{Module call graph. Shaded modules own the edge operations. In Config~A every arrow is an in-process
call; in Config~B only arrows that leave a dashed box cross the network (browse: 2, checkout: 4); in
Configs~C and~D every arrow does (browse: 3, checkout: 6, the Reservations$\to$Stock call nested).}
\label{fig:callgraph}
\end{figure}


An operation is served only by the process that hosts its owning module (Search for browse, Orders for
checkout); Figure~\ref{fig:callgraph} shows the resulting call graph. Table~\ref{tab:hops} gives the resulting number of remote calls. Browse needs two remote calls
in~B and three in~C; checkout needs four in~B and six in~C, one of them nested (Reservations to Stock).
The two operations therefore span zero to six hops with the same business logic.

\subsection{Equal resource budget}
\label{sec:budget}

A comparison of architectures is only meaningful if they run on the same resources. Every configuration
receives the same total CPU: \xo{appcpu} cores for the application tier and \xo{dbcpu} cores for the
database tier, split equally across its pods as Kubernetes requests equal to limits (Guaranteed QoS).
The split is deliberately naive. It is what a team gets by dividing a node pool evenly, and it exposes an
effect that proportional allocation would hide: in Config~C, the four modules that browse does not touch
keep half of the application budget idle. Section~\ref{sec:discussion} returns to this point.

Memory is not split equally. A Spring Boot process needs roughly 300\,MiB of resident memory at idle, so an
equal split of 4\,GiB over eight services (512\,MiB each) left the edge service open to kernel OOM kills under
overload in our calibration runs. Application pods therefore get the larger of their equal share and
1\,GiB, and memory is reported as a measured cost rather than treated as an equalised input. Istio sidecars
request 0.1 and may use up to one core; this is outside the application budget, and sidecar CPU is reported
separately as part of the mesh's cost.

The JVM settings are identical everywhere: G1 is selected explicitly (under small CPU limits the JVM would
otherwise pick the serial collector for the fine-grained services and G1 for the monolith), the heap is capped
at 50\% of the container limit, Tomcat has its default 200 worker threads, and HikariCP its default ten
connections. Two settings were changed from their defaults after calibration: Tomcat accepts at most 1{,}000
connections per process, so that an overloaded service queues work in the kernel backlog instead of buffering
it in its own memory, and keep-alive connections are never closed by the server, because the JDK HTTP client
does not retry a POST that races with a server-side close.

\subsection{Per-layer latency accounting}
\label{sec:timing}

Every response carries a W3C \texttt{Server-Timing} header~\cite{W3CServerTiming} with six durations,
aggregated over the process's own work and everything it called. For a process $s$ handling one request, a
servlet filter records the arrival time and the handler measures five intervals with
\texttt{System.nanoTime}: time before the handler runs (framework dispatch, $\mathit{oth}_s$), JSON encoding
and decoding on both sides of every hop it takes part in ($\mathit{ser}_s$), JDBC calls including connection-pool
wait ($\mathit{db}_s$), the round trip $\mathit{rtt}_c$ of every outbound call $c$, and the handler's own duration.
Application time $\mathit{app}_s$ is the handler duration minus its database time and its outbound calls. By
construction
\begin{equation}
\mathit{tot}_s = \mathit{oth}_s + \mathit{ser}_s + \mathit{app}_s + \mathit{db}_s + \sum_{c} \mathit{rtt}_c .
\label{eq:tot}
\end{equation}
The callee of call $c$ returns its own total $\mathit{tot}_c$, and the caller attributes the difference to the
network,
\begin{equation}
\mathit{net}_c = \mathit{rtt}_c - \mathit{tot}_c ,
\label{eq:net}
\end{equation}
which covers everything between the two processes: the HTTP client and server stacks, kernel networking,
kube-proxy and, in Config~D, both sidecars. Both terms of~\eqref{eq:net} are differences of timestamps taken
on one host, so no clock synchronisation is needed. This replaces the span-difference method of our earlier
protocol, which subtracted timestamps taken on different nodes and depended on NTP or PTP accuracy. Each
process adds its callees' layers to its own, so the edge process reports
$\mathit{tot} = \sum\mathit{app} + \sum\mathit{db} + \sum\mathit{ser} + \sum\mathit{net} + \sum\mathit{oth}$
over the whole call tree. The load generator's observed latency minus the edge's $\mathit{tot}$ is the
\emph{edge} layer: the client-to-edge network plus any time a request waits in the edge's accept queue.
Every millisecond of end-to-end latency is thus assigned to exactly one of six layers. Because the layers
add up exactly, the earlier protocol's reconciliation check (layers versus end-to-end within 5\%) is no longer
informative; what can be checked is that each layer behaves plausibly, for instance that $\mathit{net}$ is zero
in Config~A, which it was in every run. The probes cost about twenty \texttt{nanoTime} calls and one header of
roughly 120 bytes per hop, well under a microsecond of CPU, whereas distributed-tracing agents have been measured
to add up to 175\% latency~\cite{Nou2025}. The \texttt{calls} field counts remote calls over the call tree and
was used to confirm Table~\ref{tab:hops}.

\subsection{Functional equivalence}

Before the measurements of every deployment, the load generator replays the first 100 requests of each
workload's request sequence at 20 requests per second and keeps the response bodies. The analysis hashes
every body and compares it with the body the monolith returned for the same request, matched by sequence
number. Checkout responses
are deterministic (user, item count, total, status), so both operations can be compared byte for byte.
\xo{equiv/match} of \xo{equiv/total} bodies matched.

\section{Experimental method}
\label{sec:method}

\subsection{Environment}

The campaign reported here ran on a kind~0.31 Kubernetes~\xo{k8s} cluster on one Apple-silicon workstation
(12 cores, 24\,GiB; Docker Desktop VM with 12 vCPUs and 11.7\,GiB): one control-plane node, one node reserved
for the load generator and two nodes for the system under test, selected with node labels. All nodes share
the host's CPUs, so this environment isolates the software stack (Kubernetes networking, CFS quotas,
sidecars) but not hardware. The same manifests and scripts run unchanged on a bare-metal cluster and on
GKE, the two further environments of the protocol; Section~\ref{sec:threats} discusses what the single-host
environment can and cannot show.

\subsection{Workloads and load generation}

Load comes from Vegeta~12.13~\cite{Vegeta} running in a pod on the load-generator node. For each workload the
generator cycles through a fixed sequence of 2{,}000 requests with pseudo-random users, categories and pages,
generated from one seed, so every configuration receives the same request sequence and only the host name
differs. Two load modes answer two different questions~\cite{Schroeder2006}:

\begin{itemize}
\item \textbf{Open loop, for latency.} Requests arrive at a constant rate regardless of responses, so a slow
system cannot slow the generator down and latency is measured from each request's scheduled send time; this
avoids coordinated omission~\cite{Tene2015}. Each rate level runs for \xo{openstep}\,s. A level is
\emph{saturated} if more than 5\% of requests fail (including 10\,s client timeouts) or if throughput falls
below 90\% of the offered rate; the sweep for a configuration stops after two consecutive saturated levels.
\item \textbf{Closed loop, for capacity.} $N$ workers each send the next request as soon as the previous one
completes, with no think time, for $N \in \{1, 2, 4, \dots, 256\}$ and \xo{closedstep}\,s per level. This is
the setting in which the USL~\eqref{eq:usl} is defined, so USL fits use closed-loop data only.
\end{itemize}

Each workload in each deployment starts with a \xo{warmup}\,s closed-loop warm-up with 16 workers, which is
discarded. A closed-loop warm-up runs at each configuration's own capacity, so every JVM serves tens of thousands
of requests before measurement (JIT compilation, connection pools, caches) and none is overloaded. Calibration
runs with a short open-loop warm-up showed p99 falling as the offered rate rose, a sign of JIT compilation still
in progress during the first steps.

\subsection{Procedure}

A \emph{cell} is one configuration in one repetition. Every cell gets a fresh deployment: the namespace is
deleted and recreated, databases are re-seeded, and all JVMs start cold. Repetitions are therefore independent
samples of the whole deployment, not only of the measurement. Within each repetition the order of the four
configurations is shuffled with a fixed seed, so slow drifts of the host do not line up with one configuration.
After every step the runner records per-container CPU and memory (\texttt{kubectl top}), each container's
cgroup \texttt{cpu.stat} counters before and after the step (CPU time, CFS periods and throttled periods), and
the pods' restart counts. A step during which any pod restarted, or during which the host was suspended (detected by
comparing wall-clock time with an uptime clock that stops during sleep), is recorded as an incident and excluded from
the analysis. After a saturated step, and before each closed-loop sweep, the runner waits until a 3\,s probe at 20
requests per second completes without errors and with a p99 below 500\,ms, so that the backlog of an overloaded step
cannot leak into the next measurement; \xo{cells} cells ran \xo{reps} repetitions in \xo{hours}\,h.

\subsection{Statistical analysis}
\label{sec:stats}

The unit of replication is the cell. Requests within one run are serially dependent, so treating each request as
an independent observation would overstate precision. Per-step statistics (percentiles, mean layers,
throughput) are computed within a cell; across repetitions we report the median and, with five or more
repetitions, a 95\% bootstrap confidence interval of that median (with fewer, the range).

Tail-latency comparisons between a microservice configuration and the monolith use a two-level bootstrap in the
spirit of Kalibera and Jones~\cite{Kalibera2013}: resample repetitions with replacement, then requests within each
resampled repetition, recompute the pooled percentile for both configurations, and take the ratio. The
confidence interval of the p99 ratio is the percentile interval of the bootstrap distribution. For the H1 family
(three microservice configurations times two workloads) the interval level is Bonferroni-adjusted to
$1 - 0.05/6$. With few repetitions this procedure is more informative than a rank test on per-repetition
percentiles, whose smallest attainable two-sided $p$-value with three repetitions per group is $0.1$.

The USL is fitted by nonlinear least squares (trust-region reflective, $\lambda, \alpha, \beta \ge 0$) to every
closed-loop measurement of a configuration, one point per repetition and concurrency level. Confidence intervals
for $\alpha$, $\beta$, $N^\ast$ and the peak throughput $X_{\max}$ come from 1{,}000 bootstrap refits in which, at
every concurrency level, the measurement of a randomly drawn repetition is used. When $\beta$ is estimated as zero
the curve has no retrograde region and $N^\ast$ is undefined; $X_{\max}$ is then the asymptote $\lambda/\alpha$.

We define two crossovers and report them separately, because they need not coincide. The \emph{throughput
crossover} of configuration $X$ is the smallest concurrency at which its fitted USL curve exceeds the monolith's.
The \emph{latency crossover} is the smallest offered rate at which $X$'s median p99 is lower than the monolith's,
counting a saturated level as infinitely slow.

\subsection{Hypotheses and falsification conditions}
\label{sec:hypotheses}

The three hypotheses of the protocol are kept, with falsification conditions tied to the procedures above.

\textbf{H1 (low-load overhead).} At the lowest offered rate of each workload, the monolith has a lower p99
than every microservice configuration. H1 is falsified if the Bonferroni-adjusted interval of any ratio
$\mathrm{p99}_X/\mathrm{p99}_A$ lies entirely below one; it is supported if all six intervals lie above one.

\textbf{H2 (scaling crossover).} Under the equal budget, some microservice configuration overtakes the
monolith's throughput within the tested concurrency range, and its coherency coefficient is lower. H2 is
falsified if the monolith's throughput matches or exceeds every microservice configuration at every measured
concurrency level, or if no microservice configuration has $\beta < \beta_A$.

\textbf{H3 (mesh compounding).} Istio's overhead grows with call-graph depth and exceeds a quarter of end-to-end
latency on the six-hop path. With C and D differing only in the mesh, the overhead at a given offered rate is
$D - C$: per-hop inter-service transit, database round-trip time, and end-to-end latency. The difference isolates
the mesh only while neither configuration is limited by its CPU quota, so it is computed at the rates at which
neither saturated and the edge application container of each was throttled in less than 1\% of its CFS periods. H3 is falsified if the
mesh accounts for less than 25\% of D's end-to-end latency on checkout (six hops), or if the per-hop cost on
checkout is more than 20\% lower than on browse (three hops), both taken at the workload's lowest offered
rate, which would indicate sub-linear growth.
